\section{Problem Formulation}
\label{sec:problem}

For tracked Actor $i$, let $Z_i=(ID_i,T_{i,0:H},D_i,C_i)$ denote identity, rigid trajectory,
dimensions, and category. Build observations $B_i$ contain LiDAR endpoints and their sensor origins.
Disjoint target sweeps $Y_i$ supply surface and first-return labels.
We learn a canonical physical state $P_i=(S_i,m_i)$, where $S_i$ contains primitive centers and support
scales, and $m_i$ contains continuous evidence. The Actor state $Z_i$ is preserved during reconstruction.

For a target ray $r=(o,v,d^\star)$, a return $\widehat d$ is early when
$\widehat d<d^\star-\tau$ and a hit when $|\widehat d-d^\star|\leq\tau$.
We use $\tau=0.20$\,m. The literal point-surface operator takes the smallest positive projected depth of
any center within a $0.20$\,m lateral beam tube; an empty tube has depth $\infty$.
The categorical operator instead returns the interpolated median of the surface's depth distribution.
Both are scored against the same measured endpoint.

The learning objective combines surface reconstruction with ray agreement. Evaluation reports early rate,
hit recall, and symmetric Chamfer together: each measures a distinct consequence of changing $S_i$.
An Actor's hazard label is used only for stratified evaluation. A read-only pose places $P_i$ in world
coordinates, and a separately parameterized visual state $V_i$ supplies appearance.
